\chapter{Seeing almost Nothing}
\label{chap:almostnothing}

A picture of the supraverse runs into an awkward truth: there is almost nothing to
draw. The instanton verdict of Chapter~\ref{chap:void} is that universes are a
\emph{dilute} scatter---isolated specks in an overwhelmingly empty void---so an
honest, literal portrait is a uniform grey field with, at most, a single
infinitesimal dot. The void is the largest thing there is: nearly all of the
supraverse is three-dimensional emptiness, and the universes are tiny pockets where
a bounce has opened a fourth direction---time---inside. Visualisation here is the
art of seeing almost nothing, and of magnifying the rare somethings just enough to
make the structure legible without lying about how empty the whole is.

Before magnifying anything, it helps to see where every object in the framework
sits at once. The mass--radius plane (\Cref{fig:massradius}) is the other way to
view the whole structure: a single map on which the void seed, the black holes,
our own universe, the firstborn OGU, and the twelve fermions all fall into place
between two exclusions, with the TARDIS---a black hole read as a universe from
the inside---as its organizing idea.

\begin{figure}[p]
  \centering
  \includegraphics[width=\textwidth]{ed_massradius}
  \caption{\textbf{The Eternal Dawn mass--radius diagram.} Every object sits in the
    wedge between two exclusions: the \emph{torsion--gravity membrane} (the black-hole
    line $R=2GM/c^2$; cross it and torsion bounces out a child universe) and the
    \emph{Compton} limit (quantum uncertainty), meeting at the Planck apex---the
    smallest possible black hole. A black hole is \emph{two} points at one mass: a
    compact \emph{outside} on the membrane ($\rho\propto1/M^2$) and a cosmic
    \emph{inside}, a whole universe---the TARDIS. The blue arc is that inside: an
    enormous gap at the mountain-mass \emph{void seed}, narrowing up the line until
    inside and outside converge at the \emph{OGU} (de Sitter, $\Omega\to1$); continued
    downward (dashed) it peels off to the Compton line, where an evaporating hole's
    inside has shrunk to a single quantum---a torsiton---the Hawking end. Above the
    \emph{Hawking watershed} ($T_H=T_{\mathrm{CMB}}$, $\sim$Moon mass) a hole grows; below
    it an unfed hole evaporates, so the sub-watershed void seed is carried up only by
    the parent's feed---the magenta climb, dark matter $+$ dark energy, the ongoing
    dawn. The twelve fermions are the torsiton generations on the Compton limit (and
    tabulated, lower right), created in the genesis cascade---the dashed
    constant-density epoch lines---inside the young universe, on the same plane as the
    cosmology. Schematic; cgs, with $E=mc^2$ on the right axis. After the ``all objects''
    mass--radius diagram of Lineweaver and Patel~\citep{lineweaver2023}, reworked for
    Eternal Dawn.}
  \label{fig:massradius}
\end{figure}

\section{The challenge}

Several difficulties compound. The supraverse is a four-dimensional manifold of
non-trivial topology, and 4D does not render to a 2D page. Its structure spans
every scale from subatomic to cosmic, and it is recursively nested---universes
inside black holes inside universes. Each universe's internal time runs
orthogonally to every other's, and the whole is a static block while each pocket
within it evolves. And, above all, it is \emph{dilute}: the honest scales are so
lopsided that no linear map can show a universe and its neighbour on the same page.
Standard scientific visualisation does not handle any of this; the framework needs
its own conformal trick.

\section{The gravity-scaled conformal mapping}

The trick is to choose the one conformal rescaling that makes \emph{gravitational}
weight, not coordinate distance, set how big a thing looks. Multiply the metric by
a factor tied to the local field,
\[
  g_{\mu\nu}(x)\;\to\;\Omega^2(x)\,g_{\mu\nu}(x),
  \qquad \Omega(x)\sim\bigl(G\rho(x)\bigr)^{-1/2},
\]
the inverse square root of the local mass--energy density. High-density regions
(black-hole interiors, bounce membranes) scale \emph{up}; the low-density void
scales \emph{down}; and because the transformation is conformal it preserves causal
structure---light cones map to light cones---so what is distorted is only apparent
size, never the framework's physics. This is a cousin of the Penrose diagram, which
also uses a conformal map to fit infinite spacetime on finite paper while keeping
causal structure~\citep{penrose2010}; the difference is the choice of factor. A
Penrose diagram is shaped to fit the page; this map is shaped by \emph{gravity}, so
it foregrounds physically meaningful structure---bounces, horizons, nested
worlds---rather than coordinate artefacts.

\section{Two tiers, and why an empty picture is the honest one}
\label{sec:two-tier}

Made precise, the map magnifies a mass $M$ relative to its true size by
$\mu=(\rho/\rho_{\mathrm{ref}})^{1/3}$, and only \emph{ratios} of densities matter,
so the reference cancels. The result is two-tier, and the boundary is exactly the
identity that defines a universe (Chapter~\ref{chap:sims}, Phase~0: a universe sits
at its own Schwarzschild radius, $R_s=R_H$). An OGU is a horizon-mass black hole,
so its mean density equals the cosmic critical density---essentially the void's
dark-energy density---and therefore $\mu\simeq1$: \textbf{an OGU is not magnified
relative to its void.} At the OGU/void level, gravity-scaled coordinates simply
\emph{are} the true coordinates, and the picture is honestly near-empty in either
system. One level down it is the reverse: an astrophysical black hole is far denser
than its host, $\mu\sim10^{9}$--$10^{14}$, so the map blows these otherwise-invisible
compact objects up until the nested next-generation universes inside them become
visible (\Cref{fig:gravityscale}). The conformal factor earns its name exactly
where it must---inside a world, revealing its children---while leaving the
emptiness of the whole intact.

\begin{figure}[t]
  \centering
  \includegraphics[width=\textwidth]{gravity_scale}
  \caption{Gravity-scaling magnification $\mu=(\rho/\rho_{\mathrm{ref}})^{1/3}$.
    \emph{Left:} an OGU (horizon mass, mean density $=$ cosmic) sits at $\mu\simeq1$
    ---gravity-scaled coordinates equal true coordinates at the OGU/void level---
    while astrophysical black holes, far denser than their host, magnify by
    $10^{9}$--$10^{14}$. \emph{Right:} the black-hole mean density crosses the
    cosmic/void line precisely at the horizon mass, which is \emph{why} an OGU
    scales to unity (the $R_s=R_H$ identity). Reproducible as
    \texttt{figures/scripts/ch10\_gravity\_scale.py}.}
  \label{fig:gravityscale}
\end{figure}

\section{One number for the emptiness}

How empty is the whole? Take the ratio of the distance between neighbouring OGUs to
an OGU's own diameter. The separation is $\sim e^{I/4}$ de Sitter horizon radii and
an OGU spans $\sim2$ of them, so
\[
  \frac{\text{distance between universes}}{\text{universe diameter}}
  \;\sim\;\frac{e^{I/4}}{2}\;\sim\;10^{\,2.7\times10^{83}} ,
\]
the same whether gravity-scaled or not, since a universe and its void share a
density and neither is magnified relative to the other. The \emph{exponent} alone,
$\sim2.7\times10^{83}$, already dwarfs every number in physics---the observable
universe holds about $10^{80}$ atoms---so each universe is a lone speck whose nearest
sibling lies a gap away with more zeros than a universe has atoms. That is ``seeing
almost nothing'' made quantitative, and it is why the rendered supraverse must be
gravity-scaled and illustrative: a literal linear frame would be uniform void with a
single un-renderable dot.

\section{The honest portrait, and the wallpaper}

\Cref{fig:dilute} is the portrait the computation actually selects: a dominating,
faintly-speckled void with universes as sparse, \emph{round}, isolated islands---
gravity wells with chirality-toned membranes and faint nested children---never a
packed polygonal foam. (The foam was the legible \emph{percolating} limit, the
$I<I_{\mathrm{crit}}$ case the instanton calculation rules out; it survives only as a
contrast.) Roundness is real, not artistic licence: the vacuum has weight, so the
void occupies genuine volume and presses the islands into discs rather than cells.

The book's cover is the same picture, stripped to its recursion: a uniform void; a
balanced scatter of OGUs edged black for matter and white for antimatter; inside
each, the same-coloured first-generation children; and inside those, the dots that
are the black holes we actually watch form in our sky---the next generation down.
It is a near-empty field with richly magnified nests, which is the truthful shape of
a dilute, recursive, ever-dawning supraverse: almost nothing, lit here and there by
a world still in its dawn.

\begin{figure}[t]
  \centering
  \includegraphics[width=\textwidth]{dilute_supraverse}
  \caption{The dilute supraverse, gravity-scaled (sizes and spacings illustrative,
    not linear). The weighted void dominates; universes are sparse, \emph{round},
    isolated islands with chirality-toned membranes and faint nested children; the
    inset marks \emph{home}. At a literal linear scale the frame would be near-empty
    ---neighbours lie $\sim e^{I/4}$ horizons apart---so the void's weight keeps the
    islands round while the true gaps are astronomically larger than drawn.
    Reproducible as \texttt{figures/scripts/ch10\_dilute.py}.}
  \label{fig:dilute}
\end{figure}

\section{Targets and implementation}

Beyond the whole-supraverse portrait, the same gravity-scaled coordinates serve a
handful of specific renders: a single universe with its nested black holes (zoom
into any one and more universes appear, self-similarly); a family tree of lineages,
chirality-coded; the slow time-evolution within one bubble; and the structure of a
bounce membrane seen as the shared surface it is---one side a parent's horizon, the
other a child's first sky. All of it is ordinary Python with the gravity-scaled
transform applied before plotting, reproducible from the figure scripts alongside
every other number in this book. The visualisation track needs no simulation output
to begin (Chapter~\ref{chap:sims}); it is the most accessible way into the program,
and the quickest way to \emph{see}, however magnified, a cosmos that is almost
entirely nothing.
